% --- Archon physics formalization source begin ---
% archon:physics
% archon:covers PhyXMiniProblems/problem_phyx_mini_0449.lean
% archon:source-report reports/phyx_mini/problem_phyx_mini_0449.source.json

\chapter{Physics problem phyx\_mini\_0449}
\label{ch:PhyXMiniProblems_problem_phyx_mini_0449}

\paragraph{Problem source.}
\#\# Physical scenario

A \$400 \textbackslash{}, \textbackslash{}text\{L\}\$ tank, \$A\$, contains argon gas at \$250 \textbackslash{}, \textbackslash{}text\{kPa\}\$ and \$30 \textbackslash{}, \textasciicircum{}\{\textbackslash{}circ\}\textbackslash{}text\{C\}\$. Cylinder \$B\$, having a frictionless piston of such mass that a pressure of \$150 \textbackslash{}, \textbackslash{}text\{kPa\}\$ will float it, is initially empty. The valve is opened, and argon flows into \$B\$ and eventually reaches a uniform state of \$150 \textbackslash{}, \textbackslash{}text\{kPa\}\$ and \$30 \textbackslash{}, \textasciicircum{}\{\textbackslash{}circ\}\textbackslash{}text\{C\}\$ throughout.

\#\# Question

What is the work done by the argon?

\#\# Answer choices

- A: A: 60 kJ

- B: B: 100 kJ

- C: C: 150 kJ

- D: D: 40 kJ

\#\# Figure caption (auxiliary; use the image as primary evidence)

The image displays a diagram of a system involving gas and a piston. 

1. **Objects and Labels:**
   - There is a rectangular container on the left labeled "Argon" with the letter "A." It represents a volume of the gas argon.
   - A pipe or connection extends from the bottom of the argon container, leading to a circular symbol, possibly representing a valve.
   - This connection continues toward a vertical component on the right.

2. **Right Component:**
   - The right section is a vertical cylinder with an open top.
   - Inside this cylinder, there is a movable component labeled "B," representing a piston or similar object.
   - Above the piston, it is labeled with "P₀," indicating a pressure acting on the piston from above.
   - An arrow labeled "g" points downward on the right side, suggesting the direction of gravitational force.

3. **Relationships:**
   - The argon gas in container "A" is likely used to exert pressure or control the movement of piston "B" in the vertical cylinder.
   - The valve connects these two components, presumably controlling the flow or pressure between them.

The arrangement suggests a mechanism related to pressure and possibly used for demonstrating principles of thermodynamics or fluid dynamics.

\#\# Dataset metadata

Laws of Thermodynamics; Physical Model Grounding Reasoning, Multi-Formula Reasoning

\paragraph{Recorded answer/context.}
D: D: 40 kJ

\paragraph{Figure/image path.}
/root/proposal\_for\_physic/science-mango/phyx\_mini\_run/phyx\_data/test\_image/449.png

\paragraph{Formalization target.}
create a compiling Lean file with sorry bodies at `PhyXMiniProblems/problem\_phyx\_mini\_0449.lean`. The Lean declarations must preserve the physical quantities, dimensions or dimensional roles, figure labels, governing-law hypotheses, and final relation expressed by this problem.
Use Mathlib/Physlib names found through LeanExplore where available. If a physics API is missing, introduce faithful local abstractions rather than scalar placeholder aliases.

\begin{theorem}[Physics formalization target]
\label{thm:physics:phyx_mini_0449:target}
\lean{PhyXMiniProblems.ProblemPhyXMini0449.workDoneByArgon_matches_choice_D}
\uses{def:physics:phyx-mini-0449:phyxminiproblems-problemphyxmini0449-energyinkilojoules, def:physics:phyx-mini-0449:phyxminiproblems-problemphyxmini0449-argontankpistonsetup, def:physics:phyx-mini-0449:phyxminiproblems-problemphyxmini0449-matchesargontankpistonscenario, def:physics:phyx-mini-0449:phyxminiproblems-problemphyxmini0449-matchesproblemreadouts, def:physics:phyx-mini-0449:phyxminiproblems-problemphyxmini0449-matchessuppliedtankpistonfigure, def:physics:phyx-mini-0449:phyxminiproblems-problemphyxmini0449-hasphysicaltankpistonparameters, def:physics:phyx-mini-0449:phyxminiproblems-problemphyxmini0449-satisfiesidealgaspistonandworklaws, def:physics:phyx-mini-0449:phyxminiproblems-problemphyxmini0449-displayedworkinkilojoules}
The assigned autoformalize agent should translate this physics problem into Lean declarations in the covered file, with theorem and lemma proof bodies written as `by sorry`.
\end{theorem}
\begin{proof}
This is an autoformalization task, not a proof task. Produce statements that can later be proved by the physics prover without weakening the physical model.
\end{proof}

% --- Archon declaration topology begin ---
\section{Declaration topology}
\begin{definition}[Volume Quantity]
  \label{def:physics:phyx-mini-0449:phyxminiproblems-problemphyxmini0449-volumequantity}
  \lean{PhyXMiniProblems.ProblemPhyXMini0449.VolumeQuantity}
  A nonnegative physical volume carrying dimension length cubed
\end{definition}
\begin{proof}
  This is the declared model definition of Volume Quantity.
\end{proof}
\begin{definition}[Mass Quantity]
  \label{def:physics:phyx-mini-0449:phyxminiproblems-problemphyxmini0449-massquantity}
  \lean{PhyXMiniProblems.ProblemPhyXMini0449.MassQuantity}
  A nonnegative physical mass carrying the SI mass dimension
\end{definition}
\begin{proof}
  This is the declared model definition of Mass Quantity.
\end{proof}
\begin{definition}[Acceleration Quantity]
  \label{def:physics:phyx-mini-0449:phyxminiproblems-problemphyxmini0449-accelerationquantity}
  \lean{PhyXMiniProblems.ProblemPhyXMini0449.AccelerationQuantity}
  A nonnegative acceleration magnitude carrying dimension L T⁻²
\end{definition}
\begin{proof}
  This is the declared model definition of Acceleration Quantity.
\end{proof}
\begin{definition}[pressure In Pascals]
  \label{def:physics:phyx-mini-0449:phyxminiproblems-problemphyxmini0449-pressureinpascals}
  \lean{PhyXMiniProblems.ProblemPhyXMini0449.pressureInPascals}
  Read a physical pressure in coherent SI units (pascals)
\end{definition}
\begin{proof}
  This is the declared model definition of pressure In Pascals.
\end{proof}
\begin{definition}[pressure In Kilopascals]
  \label{def:physics:phyx-mini-0449:phyxminiproblems-problemphyxmini0449-pressureinkilopascals}
  \lean{PhyXMiniProblems.ProblemPhyXMini0449.pressureInKilopascals}
  \uses{def:physics:phyx-mini-0449:phyxminiproblems-problemphyxmini0449-pressureinpascals}
  Read a pressure in the kilopascals used in the problem statement
\end{definition}
\begin{proof}
  This is the declared model definition of pressure In Kilopascals.
\end{proof}
\begin{definition}[volume In Cubic Meters]
  \label{def:physics:phyx-mini-0449:phyxminiproblems-problemphyxmini0449-volumeincubicmeters}
  \lean{PhyXMiniProblems.ProblemPhyXMini0449.volumeInCubicMeters}
  \uses{def:physics:phyx-mini-0449:phyxminiproblems-problemphyxmini0449-volumequantity}
  Read a physical volume in coherent SI units (cubic metres)
\end{definition}
\begin{proof}
  This is the declared model definition of volume In Cubic Meters.
\end{proof}
\begin{definition}[volume In Liters]
  \label{def:physics:phyx-mini-0449:phyxminiproblems-problemphyxmini0449-volumeinliters}
  \lean{PhyXMiniProblems.ProblemPhyXMini0449.volumeInLiters}
  \uses{def:physics:phyx-mini-0449:phyxminiproblems-problemphyxmini0449-volumequantity, def:physics:phyx-mini-0449:phyxminiproblems-problemphyxmini0449-volumeincubicmeters}
  Read a physical volume in litres
\end{definition}
\begin{proof}
  This is the declared model definition of volume In Liters.
\end{proof}
\begin{definition}[energy In Joules]
  \label{def:physics:phyx-mini-0449:phyxminiproblems-problemphyxmini0449-energyinjoules}
  \lean{PhyXMiniProblems.ProblemPhyXMini0449.energyInJoules}
  Read signed dimensionful energy in coherent SI units (joules)
\end{definition}
\begin{proof}
  This is the declared model definition of energy In Joules.
\end{proof}
\begin{definition}[energy In Kilojoules]
  \label{def:physics:phyx-mini-0449:phyxminiproblems-problemphyxmini0449-energyinkilojoules}
  \lean{PhyXMiniProblems.ProblemPhyXMini0449.energyInKilojoules}
  \uses{def:physics:phyx-mini-0449:phyxminiproblems-problemphyxmini0449-energyinjoules}
  Read signed dimensionful energy in kilojoules
\end{definition}
\begin{proof}
  This is the declared model definition of energy In Kilojoules.
\end{proof}
\begin{definition}[area In Square Meters]
  \label{def:physics:phyx-mini-0449:phyxminiproblems-problemphyxmini0449-areainsquaremeters}
  \lean{PhyXMiniProblems.ProblemPhyXMini0449.areaInSquareMeters}
  Read a piston face area in square metres
\end{definition}
\begin{proof}
  This is the declared model definition of area In Square Meters.
\end{proof}
\begin{definition}[mass In Kilograms]
  \label{def:physics:phyx-mini-0449:phyxminiproblems-problemphyxmini0449-massinkilograms}
  \lean{PhyXMiniProblems.ProblemPhyXMini0449.massInKilograms}
  \uses{def:physics:phyx-mini-0449:phyxminiproblems-problemphyxmini0449-massquantity}
  Read a piston mass in kilograms
\end{definition}
\begin{proof}
  This is the declared model definition of mass In Kilograms.
\end{proof}
\begin{definition}[acceleration In Meters Per Second Squared]
  \label{def:physics:phyx-mini-0449:phyxminiproblems-problemphyxmini0449-accelerationinmeterspersecondsquared}
  \lean{PhyXMiniProblems.ProblemPhyXMini0449.accelerationInMetersPerSecondSquared}
  \uses{def:physics:phyx-mini-0449:phyxminiproblems-problemphyxmini0449-accelerationquantity}
  Read an acceleration magnitude in metres per second squared
\end{definition}
\begin{proof}
  This is the declared model definition of acceleration In Meters Per Second Squared.
\end{proof}
\begin{definition}[temperature In Kelvins]
  \label{def:physics:phyx-mini-0449:phyxminiproblems-problemphyxmini0449-temperatureinkelvins}
  \lean{PhyXMiniProblems.ProblemPhyXMini0449.temperatureInKelvins}
  Read absolute temperature in kelvins when Physlib's stored magnitude uses the given zero-preserving temperature unit
\end{definition}
\begin{proof}
  This is the declared model definition of temperature In Kelvins.
\end{proof}
\begin{definition}[temperature In Degrees Celsius]
  \label{def:physics:phyx-mini-0449:phyxminiproblems-problemphyxmini0449-temperatureindegreescelsius}
  \lean{PhyXMiniProblems.ProblemPhyXMini0449.temperatureInDegreesCelsius}
  \uses{def:physics:phyx-mini-0449:phyxminiproblems-problemphyxmini0449-temperatureinkelvins}
  Celsius readout, using the exact offset 273.15 K = 5463 / 20 K
\end{definition}
\begin{proof}
  This is the declared model definition of temperature In Degrees Celsius.
\end{proof}
\begin{definition}[Working Gas]
  \label{def:physics:phyx-mini-0449:phyxminiproblems-problemphyxmini0449-workinggas}
  \lean{PhyXMiniProblems.ProblemPhyXMini0449.WorkingGas}
  Chemical identity of the gas occupying the connected apparatus
\end{definition}
\begin{proof}
  This is the declared model definition of Working Gas.
\end{proof}
\begin{definition}[Gas Equation Of State]
  \label{def:physics:phyx-mini-0449:phyxminiproblems-problemphyxmini0449-gasequationofstate}
  \lean{PhyXMiniProblems.ProblemPhyXMini0449.GasEquationOfState}
  Equation-of-state model used for the argon at the stated conditions
\end{definition}
\begin{proof}
  This is the declared model definition of Gas Equation Of State.
\end{proof}
\begin{definition}[Tank Wall Behavior]
  \label{def:physics:phyx-mini-0449:phyxminiproblems-problemphyxmini0449-tankwallbehavior}
  \lean{PhyXMiniProblems.ProblemPhyXMini0449.TankWallBehavior}
  Mechanical behavior of tank A
\end{definition}
\begin{proof}
  This is the declared model definition of Tank Wall Behavior.
\end{proof}
\begin{definition}[Piston Friction Model]
  \label{def:physics:phyx-mini-0449:phyxminiproblems-problemphyxmini0449-pistonfrictionmodel}
  \lean{PhyXMiniProblems.ProblemPhyXMini0449.PistonFrictionModel}
  Friction model for the piston in cylinder B
\end{definition}
\begin{proof}
  This is the declared model definition of Piston Friction Model.
\end{proof}
\begin{definition}[Piston Load Regime]
  \label{def:physics:phyx-mini-0449:phyxminiproblems-problemphyxmini0449-pistonloadregime}
  \lean{PhyXMiniProblems.ProblemPhyXMini0449.PistonLoadRegime}
  External-load model encountered by the moving piston
\end{definition}
\begin{proof}
  This is the declared model definition of Piston Load Regime.
\end{proof}
\begin{definition}[Piston State]
  \label{def:physics:phyx-mini-0449:phyxminiproblems-problemphyxmini0449-pistonstate}
  \lean{PhyXMiniProblems.ProblemPhyXMini0449.PistonState}
  Piston positions relevant to the initially empty and final floating states
\end{definition}
\begin{proof}
  This is the declared model definition of Piston State.
\end{proof}
\begin{definition}[Valve Operation]
  \label{def:physics:phyx-mini-0449:phyxminiproblems-problemphyxmini0449-valveoperation}
  \lean{PhyXMiniProblems.ProblemPhyXMini0449.ValveOperation}
  The internal valve operation described in the prose
\end{definition}
\begin{proof}
  This is the declared model definition of Valve Operation.
\end{proof}
\begin{definition}[Inventory Regime]
  \label{def:physics:phyx-mini-0449:phyxminiproblems-problemphyxmini0449-inventoryregime}
  \lean{PhyXMiniProblems.ProblemPhyXMini0449.InventoryRegime}
  Whether the gas inventory can cross the outer boundary of the apparatus
\end{definition}
\begin{proof}
  This is the declared model definition of Inventory Regime.
\end{proof}
\begin{definition}[Final State Distribution]
  \label{def:physics:phyx-mini-0449:phyxminiproblems-problemphyxmini0449-finalstatedistribution}
  \lean{PhyXMiniProblems.ProblemPhyXMini0449.FinalStateDistribution}
  Spatial character of the final thermodynamic state
\end{definition}
\begin{proof}
  This is the declared model definition of Final State Distribution.
\end{proof}
\begin{definition}[Figure Symbol]
  \label{def:physics:phyx-mini-0449:phyxminiproblems-problemphyxmini0449-figuresymbol}
  \lean{PhyXMiniProblems.ProblemPhyXMini0449.FigureSymbol}
  Literal symbols visible in the supplied raster
\end{definition}
\begin{proof}
  This is the declared model definition of Figure Symbol.
\end{proof}
\begin{definition}[Figure Region]
  \label{def:physics:phyx-mini-0449:phyxminiproblems-problemphyxmini0449-figureregion}
  \lean{PhyXMiniProblems.ProblemPhyXMini0449.FigureRegion}
  Regions to which the raster's symbols are attached
\end{definition}
\begin{proof}
  This is the declared model definition of Figure Region.
\end{proof}
\begin{definition}[Vertical Direction]
  \label{def:physics:phyx-mini-0449:phyxminiproblems-problemphyxmini0449-verticaldirection}
  \lean{PhyXMiniProblems.ProblemPhyXMini0449.VerticalDirection}
  Direction of the gravity arrow in the figure
\end{definition}
\begin{proof}
  This is the declared model definition of Vertical Direction.
\end{proof}
\begin{definition}[Tank Piston Figure]
  \label{def:physics:phyx-mini-0449:phyxminiproblems-problemphyxmini0449-tankpistonfigure}
  \lean{PhyXMiniProblems.ProblemPhyXMini0449.TankPistonFigure}
  \uses{def:physics:phyx-mini-0449:phyxminiproblems-problemphyxmini0449-figuresymbol, def:physics:phyx-mini-0449:phyxminiproblems-problemphyxmini0449-figureregion, def:physics:phyx-mini-0449:phyxminiproblems-problemphyxmini0449-verticaldirection}
  Qualitative information represented by the primary figure. It records the labels Argon, A, B, P₀, and g, together with the valve connection and the vertical piston geometry. It contains no work value
\end{definition}
\begin{proof}
  This is the declared model definition of Tank Piston Figure.
\end{proof}
\begin{definition}[Argon Tank Piston Setup]
  \label{def:physics:phyx-mini-0449:phyxminiproblems-problemphyxmini0449-argontankpistonsetup}
  \lean{PhyXMiniProblems.ProblemPhyXMini0449.ArgonTankPistonSetup}
  \uses{def:physics:phyx-mini-0449:phyxminiproblems-problemphyxmini0449-volumequantity, def:physics:phyx-mini-0449:phyxminiproblems-problemphyxmini0449-massquantity, def:physics:phyx-mini-0449:phyxminiproblems-problemphyxmini0449-accelerationquantity, def:physics:phyx-mini-0449:phyxminiproblems-problemphyxmini0449-workinggas, def:physics:phyx-mini-0449:phyxminiproblems-problemphyxmini0449-gasequationofstate, def:physics:phyx-mini-0449:phyxminiproblems-problemphyxmini0449-tankwallbehavior, def:physics:phyx-mini-0449:phyxminiproblems-problemphyxmini0449-pistonfrictionmodel, def:physics:phyx-mini-0449:phyxminiproblems-problemphyxmini0449-pistonloadregime, def:physics:phyx-mini-0449:phyxminiproblems-problemphyxmini0449-pistonstate, def:physics:phyx-mini-0449:phyxminiproblems-problemphyxmini0449-valveoperation, def:physics:phyx-mini-0449:phyxminiproblems-problemphyxmini0449-inventoryregime, def:physics:phyx-mini-0449:phyxminiproblems-problemphyxmini0449-finalstatedistribution, def:physics:phyx-mini-0449:phyxminiproblems-problemphyxmini0449-tankpistonfigure}
  Independent physical quantities and qualitative attributes of the transfer. The work observable and final cylinder volume are fields, not definitions in terms of the requested answer; their governing relations are assumptions below
\end{definition}
\begin{proof}
  This is the declared model definition of Argon Tank Piston Setup.
\end{proof}
\begin{definition}[Matches Argon Tank Piston Scenario]
  \label{def:physics:phyx-mini-0449:phyxminiproblems-problemphyxmini0449-matchesargontankpistonscenario}
  \lean{PhyXMiniProblems.ProblemPhyXMini0449.MatchesArgonTankPistonScenario}
  \uses{def:physics:phyx-mini-0449:phyxminiproblems-problemphyxmini0449-volumeincubicmeters, def:physics:phyx-mini-0449:phyxminiproblems-problemphyxmini0449-argontankpistonsetup}
  Qualitative idealizations stated by the physical scenario
\end{definition}
\begin{proof}
  This is the declared model definition of Matches Argon Tank Piston Scenario.
\end{proof}
\begin{definition}[Matches Problem Readouts]
  \label{def:physics:phyx-mini-0449:phyxminiproblems-problemphyxmini0449-matchesproblemreadouts}
  \lean{PhyXMiniProblems.ProblemPhyXMini0449.MatchesProblemReadouts}
  \uses{def:physics:phyx-mini-0449:phyxminiproblems-problemphyxmini0449-pressureinkilopascals, def:physics:phyx-mini-0449:phyxminiproblems-problemphyxmini0449-volumeinliters, def:physics:phyx-mini-0449:phyxminiproblems-problemphyxmini0449-temperatureindegreescelsius, def:physics:phyx-mini-0449:phyxminiproblems-problemphyxmini0449-argontankpistonsetup}
  Numerical readouts supplied by the prose. The two 30 °C facts make the transfer isothermal, while the two independent 150 kPa facts describe the piston calibration and the eventual uniform gas pressure
\end{definition}
\begin{proof}
  This is the declared model definition of Matches Problem Readouts.
\end{proof}
\begin{definition}[Matches Supplied Tank Piston Figure]
  \label{def:physics:phyx-mini-0449:phyxminiproblems-problemphyxmini0449-matchessuppliedtankpistonfigure}
  \lean{PhyXMiniProblems.ProblemPhyXMini0449.MatchesSuppliedTankPistonFigure}
  \uses{def:physics:phyx-mini-0449:phyxminiproblems-problemphyxmini0449-argontankpistonsetup}
  Primary-image evidence, separated from the numerical prose readouts
\end{definition}
\begin{proof}
  This is the declared model definition of Matches Supplied Tank Piston Figure.
\end{proof}
\begin{definition}[Has Physical Tank Piston Parameters]
  \label{def:physics:phyx-mini-0449:phyxminiproblems-problemphyxmini0449-hasphysicaltankpistonparameters}
  \lean{PhyXMiniProblems.ProblemPhyXMini0449.HasPhysicalTankPistonParameters}
  \uses{def:physics:phyx-mini-0449:phyxminiproblems-problemphyxmini0449-pressureinpascals, def:physics:phyx-mini-0449:phyxminiproblems-problemphyxmini0449-volumeincubicmeters, def:physics:phyx-mini-0449:phyxminiproblems-problemphyxmini0449-areainsquaremeters, def:physics:phyx-mini-0449:phyxminiproblems-problemphyxmini0449-massinkilograms, def:physics:phyx-mini-0449:phyxminiproblems-problemphyxmini0449-accelerationinmeterspersecondsquared, def:physics:phyx-mini-0449:phyxminiproblems-problemphyxmini0449-temperatureinkelvins, def:physics:phyx-mini-0449:phyxminiproblems-problemphyxmini0449-argontankpistonsetup}
  Positivity and nondegeneracy conditions for the physical apparatus
\end{definition}
\begin{proof}
  This is the declared model definition of Has Physical Tank Piston Parameters.
\end{proof}
\begin{definition}[Satisfies Ideal Gas Piston And Work Laws]
  \label{def:physics:phyx-mini-0449:phyxminiproblems-problemphyxmini0449-satisfiesidealgaspistonandworklaws}
  \lean{PhyXMiniProblems.ProblemPhyXMini0449.SatisfiesIdealGasPistonAndWorkLaws}
  \uses{def:physics:phyx-mini-0449:phyxminiproblems-problemphyxmini0449-pressureinpascals, def:physics:phyx-mini-0449:phyxminiproblems-problemphyxmini0449-volumeincubicmeters, def:physics:phyx-mini-0449:phyxminiproblems-problemphyxmini0449-energyinjoules, def:physics:phyx-mini-0449:phyxminiproblems-problemphyxmini0449-areainsquaremeters, def:physics:phyx-mini-0449:phyxminiproblems-problemphyxmini0449-massinkilograms, def:physics:phyx-mini-0449:phyxminiproblems-problemphyxmini0449-accelerationinmeterspersecondsquared, def:physics:phyx-mini-0449:phyxminiproblems-problemphyxmini0449-temperatureinkelvins, def:physics:phyx-mini-0449:phyxminiproblems-problemphyxmini0449-argontankpistonsetup}
  The governing physical model consists of four relations: fixed-inventory ideal-gas behavior gives the isothermal pV invariant for the total argon volume in A and B; the calibrated floating pressure balances P₀ and the piston weight; a floating frictionless piston is in pressure equilibrium with the gas; and work done by the argon against the constant load is P\_float ΔV\_B. These relations contain neither the derived final volume 4/15 m³ nor the requested work 40 kJ
\end{definition}
\begin{proof}
  This is the declared model definition of Satisfies Ideal Gas Piston And Work Laws.
\end{proof}
\begin{lemma}[final Cylinder BVolume eq four fifteenths cubic Meters]
  \label{lem:physics:phyx-mini-0449:phyxminiproblems-problemphyxmini0449-finalcylinderbvolume-eq-four-fifteenths-cubicmeters}
  \lean{PhyXMiniProblems.ProblemPhyXMini0449.finalCylinderBVolume_eq_four_fifteenths_cubicMeters}
  \uses{def:physics:phyx-mini-0449:phyxminiproblems-problemphyxmini0449-volumeincubicmeters, def:physics:phyx-mini-0449:phyxminiproblems-problemphyxmini0449-argontankpistonsetup, def:physics:phyx-mini-0449:phyxminiproblems-problemphyxmini0449-matchesargontankpistonscenario, def:physics:phyx-mini-0449:phyxminiproblems-problemphyxmini0449-matchesproblemreadouts, def:physics:phyx-mini-0449:phyxminiproblems-problemphyxmini0449-hasphysicaltankpistonparameters, def:physics:phyx-mini-0449:phyxminiproblems-problemphyxmini0449-satisfiesidealgaspistonandworklaws}
  The final volume occupied by argon in cylinder B is 4/15 m³. This is a derived ideal-gas consequence, not a premise of the work theorem
\end{lemma}
\begin{proof}
  The conclusion follows from the governing relations and figure readouts named in the statement of final Cylinder BVolume eq four fifteenths cubic Meters.
\end{proof}
\begin{theorem}[work Done By Argon eq forty kilojoules]
  \label{thm:physics:phyx-mini-0449:phyxminiproblems-problemphyxmini0449-workdonebyargon-eq-forty-kilojoules}
  \lean{PhyXMiniProblems.ProblemPhyXMini0449.workDoneByArgon_eq_forty_kilojoules}
  \uses{def:physics:phyx-mini-0449:phyxminiproblems-problemphyxmini0449-energyinkilojoules, def:physics:phyx-mini-0449:phyxminiproblems-problemphyxmini0449-argontankpistonsetup, def:physics:phyx-mini-0449:phyxminiproblems-problemphyxmini0449-matchesargontankpistonscenario, def:physics:phyx-mini-0449:phyxminiproblems-problemphyxmini0449-matchesproblemreadouts, def:physics:phyx-mini-0449:phyxminiproblems-problemphyxmini0449-matchessuppliedtankpistonfigure, def:physics:phyx-mini-0449:phyxminiproblems-problemphyxmini0449-hasphysicaltankpistonparameters, def:physics:phyx-mini-0449:phyxminiproblems-problemphyxmini0449-satisfiesidealgaspistonandworklaws}
  The argon does positive boundary work of exactly 40 kJ while lifting the piston. This is the substantive answer requested by the problem
\end{theorem}
\begin{proof}
  The conclusion follows from the governing relations and figure readouts named in the statement of work Done By Argon eq forty kilojoules.
\end{proof}
\begin{definition}[Answer Choice]
  \label{def:physics:phyx-mini-0449:phyxminiproblems-problemphyxmini0449-answerchoice}
  \lean{PhyXMiniProblems.ProblemPhyXMini0449.AnswerChoice}
  Labels printed beside the four candidate answers
\end{definition}
\begin{proof}
  This is the declared model definition of Answer Choice.
\end{proof}
\begin{definition}[displayed Work In Kilojoules]
  \label{def:physics:phyx-mini-0449:phyxminiproblems-problemphyxmini0449-displayedworkinkilojoules}
  \lean{PhyXMiniProblems.ProblemPhyXMini0449.displayedWorkInKilojoules}
  \uses{def:physics:phyx-mini-0449:phyxminiproblems-problemphyxmini0449-answerchoice}
  Displayed candidate work values, interpreted in kilojoules
\end{definition}
\begin{proof}
  This is the declared model definition of displayed Work In Kilojoules.
\end{proof}
\begin{definition}[recorded Dataset Answer]
  \label{def:physics:phyx-mini-0449:phyxminiproblems-problemphyxmini0449-recordeddatasetanswer}
  \lean{PhyXMiniProblems.ProblemPhyXMini0449.recordedDatasetAnswer}
  \uses{def:physics:phyx-mini-0449:phyxminiproblems-problemphyxmini0449-answerchoice}
  The source dataset records answer label D; this is metadata only
\end{definition}
\begin{proof}
  This is the declared model definition of recorded Dataset Answer.
\end{proof}
% --- Archon declaration topology end ---
% --- Archon physics formalization source end ---
